% Preamble
\documentclass[11pt, aspectratio=169]{beamer}
	% Packages
	\usepackage[english]{babel}
	\usepackage{lipsum}
	\usepackage{mathptmx}
	\usepackage{xcolor}
	% Themes
	\usefonttheme[onlymath]{serif}
	\usetheme[nofirafonts]{focus}
	\renewcommand{\thempfootnote}{\arabic{mpfootnote}}
	% Cover
	\title{Microeconometrics Using Stata}
	\subtitle{Nonlinear optimization methods}
	\author{
		Jhon R. Ordoñez
		\footnote{\label{1} National University of San Cristóbal de Huamanga}
		\footnote{\label{2} Faculty of Economic, Administrative and Accounting Sciences}
		\footnote{\label{3} Professional School of Economics}
	}
	\date{\today}
% Body
\begin{document}
	% Cover
	\begin{frame}
		\maketitle
	\end{frame}
	% Outline
	\begin{frame}[t]
		\frametitle{Outline}
		\tableofcontents
	\end{frame}
	% Section 1
	\section{Exercise 1}
		% Slide 1
		\begin{frame}
			\frametitle{Exercise 1}
				\begin{block}{Exercise 1}
					Consider estimation of the logit model covered in \textcolor{red!50}{\textbf{section 10.5}} and
					\textcolor{red!50}{\textbf{chapter 17}}. Then $Q(\mathbf{\beta}) = \sum_{i}\left\{y_{i} \ln{\Lambda_{i}} + \left(1 - y_{i}\right)\Lambda_{i}\right\}$, where
					$\Lambda_{i} = \Lambda\left(\mathbf{x}_{i}'\mathbf{\beta}\right) = exp\left(\mathbf{x}_{i}'\mathbf{\beta}\right)/ \left\{1 + \mathbf{x}_{i}'\mathbf{\beta}\right\}$. 
					Show that $g\left(\mathbf{\beta}\right) = \sum_{i}\left(y_{i} - \Lambda_{i}\right)\mathbf{x}_{i}$
					and $H\left(\mathbf{\beta}\right) = \sum_{i} - \Lambda_{i}\left(1 - \Lambda_{i}\right)\mathbf{x}_{i}\mathbf{x}_{i}'$. 
					\textcolor{magenta!87}{Hint: \textbf{$\partial \Lambda \left(z\right) / \partial z = \Lambda\left(z\right)\left\{1 - \Lambda\left(z\right)\right\}$}}.
					Use the data on \textcolor{blue}{\textbf{docvis}} to generate the binary variable \textcolor{blue}{\textbf{d$\_$dv}} for whether there are
					any doctor visits. Using just 2002 data, as in this chapter, use \textcolor{blue}{\textbf{logit}} to
					perform logistic regression of the binary variable \textcolor{blue}{\textbf{d$\_$dv}} on \textcolor{blue}{\textbf{private}}, \textcolor{blue}{\textbf{chronic}},
					\textcolor{blue}{\textbf{female}}, \textcolor{blue}{\textbf{income}}, and an intercept. Obtain robust estimates of the standard
					errors. You should find that the coefficient of \textcolor{blue}{\textbf{private}}, for example, equals $1.27266$
					, with a robust standard error of $0.0896928$.
				\end{block}
		\end{frame}
	% Section 2
	\section{Exercise 2}
		% Slide 1
		\begin{frame}
			\frametitle{Exercise 2}
			\begin{block}{Exercise 2}
				Adapt the code of \textcolor{red!50}{\textbf{section 16.2.3}} to fit the logit model
				of exercise $1$ using \textbf{NR} iterations coded in Mata. 
				\textcolor{magenta!87}{\textbf{Hint:} In defining an $n\times 1$ column vector with entries $\Lambda_{i}$ , you may find it helpful
				to use the fact that $J(n,1,1)$ creates an $n \times 1$ vector of $1$s}.
			\end{block}
		\end{frame}
	% References
	\begin{frame}[t,allowframebreaks]
		\frametitle{References}
		\bibliography{biblio.bib}
		% Style
		\bibliographystyle{apalike}
		% Authors
		\nocite{cameron2022-I}
		\nocite{cameron2022-II}
	\end{frame}
	% Cover
	\begin{frame}
		\maketitle
	\end{frame}
\end{document}